% --- Archon physics formalization source begin ---
% archon:physics
% archon:covers IPhO2026Problems/problem_ipho_2026_t3_a3.lean
% archon:source-report reports/ipho_2026/problem_ipho_2026_t3_a3.source.json
% archon:problem-id ipho_2026_t3
% archon:part-id T3-A3
% archon:previous-part ipho_2026_t3_a2
% archon:previous-part-policy derive_inline_from_problem_only_material

\chapter{Ipho Physics problem ipho\_2026\_t3\_a3}
\label{ch:IPhO2026Problems_problem_ipho_2026_t3_a3}

\paragraph{Problem source.}
\#\# Physical scenario

A homogeneous isotropic paramagnetic torus has mean radius R, inner radius r
with r << R, volume V, and cross-sectional area A.  An insulated conducting
wire is wound densely around it with N turns and instantaneous current I.
Fields H and B and magnetization M are approximately uniform in the torus.
Use B = mu\_0*H + mu\_0*M, Ampere's law, and the sign convention that work and
heat entering the paramagnetic torus are positive.

\#\# Current subquestion T3-A3

Subtract the vacuum-core contribution and write the work dW done on the paramagnetic material.

\paragraph{Shared problem context.}
A homogeneous isotropic paramagnetic torus has mean radius R, inner radius r
with r << R, volume V, and cross-sectional area A.  An insulated conducting
wire is wound densely around it with N turns and instantaneous current I.
Fields H and B and magnetization M are approximately uniform in the torus.
Use B = mu\_0*H + mu\_0*M, Ampere's law, and the sign convention that work and
heat entering the paramagnetic torus are positive.

\paragraph{Current subquestion.}
Subtract the vacuum-core contribution and write the work dW done on the paramagnetic material.

\paragraph{Figure/image paths.}
\begin{itemize}
\item ipho\_2026\_source/image/T3\_page-2.png
\item ipho\_2026\_source/image/T3\_page-1.png
\end{itemize}

\paragraph{Reusable previous-part conclusions.}
\begin{itemize}
\item Source T3-A2. Question: Find the work dW\_emf performed by the external voltage source when B changes by dB. Policy: derive\_inline\_from\_problem\_only\_material
\end{itemize}

\paragraph{Formalization target.}
create a compiling Lean file with sorry bodies at `IPhO2026Problems/problem\_ipho\_2026\_t3\_a3.lean`.
The Lean declarations must preserve the physical quantities, dimensions or dimensional roles, figure labels, governing-law hypotheses, and final relation expressed by this problem.
Use Mathlib/Physlib names found through LeanExplore where available. If a domain API is missing, introduce faithful local abstractions rather than scalar placeholder aliases.
This is an answer-blind solve-phase task. Derive a candidate result only from the problem-side evidence above; do not assume a target value, select a tolerance around a desired value, or encode a desired result into a candidate domain.
For numerical questions, define the raw end-to-end quantity and apply an explicit source-derived rounding rule. For identification questions, characterize or prove uniqueness before recording the derived candidate.

\begin{theorem}[Ipho Physics formalization target]
\label{thm:physics:ipho_2026_t3_a3:target}
The assigned autoformalize agent should translate this IPhO physics problem into Lean declarations in the covered file, with theorem and lemma proof bodies written as `by sorry`.
\end{theorem}
\begin{proof}
This is an autoformalization task, not a proof task. Produce faithful statements that can later be proved without weakening the source contract.
\end{proof}
% --- Archon physics formalization source end ---
